%@ pages 333-335
%@ status translated
% Chapter I, section 4: Additive categories (pp. 333-335).  The section opens
% in the middle of p. 333, just below the last paragraph of section 3, and
% runs past p. 334: Proposition 3 (additive functors) and the paragraphs
% following it stand at the top of p. 335, above the heading of section 5.
% Transcribed from the page images in reference/source.pdf (pages 333--335),
% not from the OCR text layer, which loses the subscripts of the canonical
% morphisms and misreads the bold Sigma and Pi as ordinary capitals.
%
% Emendations (each with an \ednote at the point):
%   * The neutral element is 0_{M,N} (p. 334 prints o_{N,M}, transposed).
%   * c(M) is inverse to 1_M in Hom(M,M) (p. 334 prints Hom(M,N)).
% The zero object and zero morphisms are rendered 0 throughout, per the
% house convention.
%
% The original prints the section heading as
%   "4. Catégories additives (d'après Grothendieck)"
% The attribution is carried by the heading rather than by the running text.

\section{Additive categories (after Grothendieck)}
--- Let $\Cat$ be a category satisfying the following axioms:

\begin{description}
\item[$\mathbf{C}$ Ad 1] There exists a zero object $0$.

\item[$\mathbf{C}$ Ad 2] For every pair $(M,N)$ of objects of $\Cat$, the
direct product of $M$ and $N$, and likewise the direct sum of $M$ and $N$,
exist.
\end{description}

\begin{lemma}
If $M$ and $N$ are two objects of $\Cat$, there exists one and only one
morphism $h(M,N)\colon M\boldsymbol{\Sigma}N\to M\boldsymbol{\Pi}N$ such
that we have
\[
\begin{gathered}
  q_M\circ h(M,N)\circ j_M = 1_M, \qquad
  q'_N\circ h(M,N)\circ j_M = 0,\\
  q'_N\circ h(M,N)\circ j'_N = 1_N, \qquad
  q_M\circ h(M,N)\circ j'_N = 0.
\end{gathered}
\]
\end{lemma}

If we denote by $\boldsymbol{\Pi}$ (resp.\ $\boldsymbol{\Sigma}$) the functor
\[
  (M,N)\leadsto M\boldsymbol{\Pi}N
  \quad[\text{resp.}\ (M,N)\leadsto M\boldsymbol{\Sigma}N]
\]
of the product category $\Cat\boldsymbol{\Pi}\Cat$ into $\Cat$, Lemma~1 can be
completed by the following:

\begin{lemma}
The morphisms $h(M,N)$ define a natural transformation of the functor
$\boldsymbol{\Sigma}$ into the functor $\boldsymbol{\Pi}$.
\end{lemma}

Let us now suppose that $\Cat$ satisfies one axiom more:

\begin{description}
\item[$\mathbf{C}$ Ad 3] For every pair $(M,N)$ of objects of $\Cat$, the
morphism $h(M,N)$ of $M\boldsymbol{\Sigma}N$ into
$M\boldsymbol{\Pi}N$ is an isomorphism.
\end{description}

Under these conditions, we may suppose that $M\boldsymbol{\Sigma}N$ has been
chosen equal to $M\boldsymbol{\Pi}N$ and such that $h(M,N)$ is the identity
morphism of $M\boldsymbol{\Pi}N$. We then write $M\oplus N$ in place of
$M\boldsymbol{\Sigma}N$ or $M\boldsymbol{\Pi}N$. It is easy to prove that the
identification of the direct sum with the direct product is compatible with
``the interchange of the factors $M$ and $N$'' and with the canonical
isomorphisms of $(M\boldsymbol{\Sigma}N)\boldsymbol{\Sigma}P$ onto
$M\boldsymbol{\Sigma}(N\boldsymbol{\Sigma}P)$ and of
$(M\boldsymbol{\Pi}N)\boldsymbol{\Pi}P$ onto
$M\boldsymbol{\Pi}(N\boldsymbol{\Pi}P)$.

If $M$ is an object of $\Cat$, we denote by $\Delta_M$ (resp.\
$\boldsymbol{\Sigma}_M$) the morphism, called \emph{diagonal}, of $M$ into
$M\oplus M$ (resp.\ of $M\oplus M$ into $M$) which is defined by the
equalities $q_M\circ\Delta_M = 1_M$, $q'_M\circ\Delta_M = 1_M$ (resp.\
$\boldsymbol{\Sigma}_M\circ j_M = 1_M$,
$\boldsymbol{\Sigma}_M\circ j'_M = 1_M$). If $f$ and $g$ are two morphisms
with the same source $M$ and the same target $N$, we denote by $f+g$ the
composite of the following morphisms:
\[
  M \xrightarrow{\Delta_M} M\oplus M
    \xrightarrow{f\oplus g} N\oplus N
    \xrightarrow{\boldsymbol{\Sigma}_M} N.
\]

\begin{lemma}
The internal law $(f,g)\mapsto f+g$ is commutative and associative. The
morphism $0_{M,N}$ is a neutral element for this law.\ednote{The original
prints $o_{N,M}$, with the indices transposed; the neutral element of
$\operatorname{Hom}(M,N)$ is the zero morphism $M\to N$. We print
$0_{M,N}$.}
\end{lemma}

\begin{lemma}
Let $f,f'\colon M\to N$ and $g,g'\colon N\to P$ be morphisms of $\Cat$. We
have the equalities $g\circ(f+f') = g\circ f + g\circ f'$ and
$(g+g')\circ f = g\circ f + g'\circ f$.
\end{lemma}

There exist categories satisfying the axioms $\mathbf{C}$ Ad 1, 2 and 3 and
for which the monoids $\operatorname{Hom}(M,N)$ are not abelian groups. This
cannot happen if the category $\Cat$ satisfies the axiom $\mathbf{C}$ Ad 4:

\begin{description}
\item[$\mathbf{C}$ Ad 4] For every object $M$, there exists a morphism
$c(M)\colon M\to M$ such that the following diagram is commutative:
\end{description}

\[
\begin{tikzcd}
M \arrow[r, "0"] \arrow[d, "\Delta_M"'] & M \\
M \oplus M \arrow[r, "1_M \oplus c(M)"] & M \oplus M \arrow[u, "\boldsymbol{\Sigma}_M"']
\end{tikzcd}
\]

The commutativity of the diagram means that $c(M)$ is an inverse of $1_M$ in
$\operatorname{Hom}(M,M)$.\ednote{The original prints
$\operatorname{Hom}(M,N)$; $c(M)$ and $1_M$ are endomorphisms of $M$. We
print $\operatorname{Hom}(M,M)$.}
It follows that a morphism $f\colon M\to N$ has for its inverse the composite
\[
  c(N)\circ f = f\circ c(M).
\]

\begin{proposition}
If $\Cat$ is a category, the two following assertions are equivalent:

\begin{description}
\item[\normalfont a.] The category $\Cat$ satisfies the axioms $\mathbf{C}$
Ad 1, 2, 3 and 4.

\item[\normalfont b.] The category $\Cat$ satisfies $\mathbf{C}$ Ad 1 and
either of the two assertions of $\mathbf{C}$ Ad 2; moreover, the sets
$\operatorname{Hom}(M,N)$ can be given a structure of abelian group in such a
way that the laws of composition are bilinear maps.
\end{description}
\end{proposition}

The implication $(a)\Rightarrow(b)$ follows from the preceding lemmas. The
implication $(b)\Rightarrow(a)$ is classical; it is proved in \cite{grothendieck-1}.

We call a category \emph{additive} if it satisfies the axioms $\mathbf{C}$
Ad~1, 2, 3 and 4.

\begin{proposition}\label{prop:additive-functor}
Let $\Cat$ and $\mathbf{D}$ be two additive categories and let $F$ be a
functor of $\Cat$ into $\mathbf{D}$. The following assertions are equivalent:

\begin{description}
\item[\normalfont a.] For every pair $(M,N)$ of objects of $\Cat$, the map
$F(M,N)$ of $\operatorname{Hom}_{\Cat}(M,N)$ into
$\operatorname{Hom}_{\mathbf{D}}(FM,FN)$ is linear.

\item[\normalfont b.] For every pair $(M,N)$ of objects of $\Cat$, the triple
\[
  \bigl(F(M \oplus N),\, Fj_M,\, Fj'_N\bigr)
\]
is a direct sum of $FM$ and $FN$.
\end{description}
\end{proposition}

If the equivalent conditions of Proposition~\ref{prop:additive-functor} are
satisfied, we say that $F$ is an \emph{additive functor}. Condition (b) is
expressed by saying that an additive functor commutes with direct sums. Of
course, an additive functor commutes with direct products as well.

Unless expressly stated to the contrary, throughout what follows, the
\emph{functors from one additive category into another are assumed to be
additive}. Let, for example, $\Cat$ be an additive category. We have seen
that, for every pair $(M,N)$ of objects of $\Cat$, the set
$\operatorname{Hom}_{\Cat}(M,N)$ is endowed with a ``natural'' structure of
abelian group. We denote this abelian group again by
$\operatorname{Hom}_{\Cat}(M,N)$. Lemma~4 then shows that the functor
$Y\leadsto \operatorname{Hom}_{\Cat}(M,Y)$ of $\Cat$ into $\Ab$ is additive
for every $M$.

We shall also observe that the dual category of an additive category $\Cat$
is additive. Thus $X\leadsto \operatorname{Hom}_{\Cat}(X,N)$ is an additive
functor of $\Cat^{\circ}$ into $\Ab$.
